\documentclass[11pt]{article}
\usepackage[a4paper,margin=27mm]{geometry}
\usepackage{amsmath,amssymb,amsthm,mathtools}
\usepackage[T1]{fontenc}
\usepackage[utf8]{inputenc}
\usepackage{lmodern}
\usepackage{microtype}
\usepackage{xcolor}
\usepackage{hyperref}
\hypersetup{colorlinks=true,linkcolor=blue!55!black,citecolor=blue!55!black,urlcolor=blue!55!black}
\setlength{\parindent}{0pt}
\setlength{\parskip}{5pt plus 1pt}
\newtheorem{theorem}{Theorem}[section]
\newtheorem{lemma}[theorem]{Lemma}
\newtheorem{proposition}[theorem]{Proposition}
\theoremstyle{remark}
\newtheorem{remark}[theorem]{Remark}
\newcommand{\per}{\operatorname{per}}
\newcommand{\one}{\mathbf 1}
\newcommand{\RR}{\mathbb R}
\newcommand{\QQ}{\mathbb Q}

\title{Dittert's conjecture in dimension 8\\
\large An exact computer-assisted boundary-exclusion proof}
\author{Pedro Paulo Marques do Nascimento}
\date{23 July 2026}

\begin{document}
\maketitle

\begin{center}
\fbox{\parbox{0.90\textwidth}{\textbf{Status.} This is an exact, reproducible
working proof. It has not yet undergone independent peer review. The finite
dimension-specific calculations are checked by two separately implemented
standard-library verifiers.}}

\small Licensed under \href{https://creativecommons.org/licenses/by/4.0/}{CC BY 4.0}.
\end{center}

\begin{abstract}
For a nonnegative $n\times n$ matrix $A$ whose entries sum to $n$, let $r_i,c_j$
be its row and column sums and put
\[
 \Phi(A)=\prod_i r_i+\prod_j c_j-\per A.
\]
We prove Dittert's conjecture in dimension $n=8$.  A maximal Li scaling
reduces the boundary problem to whole, marginal, and proper tight cuts.
Most proper cuts are excluded by exact staircase-face bounds; the six
one-dimensional-core cuts use sharper dimension-specific line-sum deficit
estimates.  For a marginal cut, stationarity forces a large column containing
the minimum row.  Cofactor nonnegativity bounds every other row meeting this
column, while the zeroes in that column supply an additional staircase-face
permanent bound.  Six exact degree-65 Bernstein certificates exclude every
possible nontrivial column support.  No four-independent-zero theorem is used.
\end{abstract}

\section{Statement and published inputs}

Let
\[
 K_n=\left\{A=(a_{ij})\in\RR_{\ge0}^{n\times n}:
       \sum_{i,j}a_{ij}=n\right\},
 \qquad
 \gamma_n=\frac{n!}{n^n}.
\]
For $A\in K_n$, write
\[
 R=\prod_{i=1}^n r_i,\qquad
 C=\prod_{j=1}^n c_j,\qquad
 P=\per A.
\]
Thus $\Phi(A)=R+C-P$ and $\Phi(J_n/n)=2-\gamma_n$.

\begin{theorem}\label{thm:main}
For every $A\in K_8$,
\[
 \Phi(A)\le 2-\frac{8!}{8^8},
\]
with equality if and only if $A=J_8/8$.
\end{theorem}

We use the following published results.
\begin{enumerate}
\item Hwang proved that the only possible positive $\Phi$-maximizer in $K_n$
is $J_n/n$ \cite{Hwang1986}.
\item Cheon and Wanless proved that a partly decomposable matrix cannot
maximize $\Phi$, and that a maximizer cannot have its complete zero set equal,
up to row and column permutations, to one proper rectangular block
\cite[Theorems 3.2 and 3.7]{CheonWanless2012}.
\item Li's cut characterization says that a nonnegative $n\times n$ matrix
$S$ is doubly superstochastic if and only if
\[
 \sum_{i\in I,j\in J}s_{ij}\ge |I|+|J|-n
\]
for all $I,J\subseteq[n]$ \cite{Li1990}; see also
\cite[Lemma 2.2]{CheonWanless2012}.
\item On a staircase face of the Birkhoff polytope, the barycenter minimizes
the permanent \cite{Hwang1985}.
\item Minc's averaging theorem for permanent minimizers with repeated support
rows and columns gives the block form used below \cite[Theorem 1]{Minc1984};
see also \cite[Section 2]{PulaSongWanless2011}.
\item The van der Waerden permanent theorem gives $\per B\ge\gamma_n$ for
every doubly stochastic $B$, with equality only at $J_n/n$
\cite{Egorychev1981,Falikman1981}.
\end{enumerate}
The support-stationarity equation and the large-column support argument used
below are proved directly; they are not additional external assumptions.

\section{Boundary structure and the shared deficit}

Because $K_n$ is compact, $\Phi$ has a global maximizer.  By Hwang's theorem,
a nonuniform maximizer must lie on the boundary.

\begin{lemma}[Two independent zeroes]\label{lem:independent-zeros}
Every boundary $\Phi$-maximizer has two zero entries in distinct rows and
distinct columns.
\end{lemma}
\begin{proof}
Suppose a nonempty zero set contains no such pair.  Fix a zero at $(i,j)$.
Every other zero shares row $i$ or column $j$.  If there were both a zero
$(i,j')$, $j'\ne j$, and a zero $(i',j)$, $i'\ne i$, those two zeroes would be
independent.  Hence all zeroes lie in one row or all lie in one column.

If that row or column is entirely zero, the matrix is partly decomposable.
Otherwise its complete zero set is one proper rectangular block.  Both cases
are excluded for a maximizer by Cheon--Wanless.
\end{proof}

Assume for contradiction that a boundary maximizer $A$ satisfies
$\Phi(A)\ge2-\gamma_n$.  Since the row sums and column sums each have
arithmetic mean $1$, AM--GM gives $R,C\le1$.  Consequently $P\le\gamma_n$.
Write
\[
 P=\gamma_n-\delta,\qquad 0\le\delta\le\gamma_n,
\]
and put $R=1-\rho$, $C=1-\sigma$.  The assumed lower bound gives
\begin{equation}\label{eq:shared-deficit}
 \rho+\sigma\le\delta.
\end{equation}
In particular,
\begin{equation}\label{eq:product-lower}
 R\ge1-\delta,\qquad C\ge1-\delta,
 \qquad RC\ge1-\delta.
\end{equation}

\section{The maximal Li scaling}

For $I,J\subseteq[n]$, let $s(A[I,J])$ be the sum of the indicated submatrix,
and define
\begin{equation}\label{eq:lambda}
 \lambda=
 \min_{\substack{I,J\subseteq[n]\\k:=|I|+|J|-n>0}}
 \frac{s(A[I,J])}{k}.
\end{equation}
A maximizer is fully indecomposable by the Cheon--Wanless exclusion.  Hence
every numerator in \eqref{eq:lambda} is positive, so $\lambda>0$; taking
$I=J=[n]$ shows $\lambda\le1$.

Li's theorem says that $A/\lambda$ is doubly superstochastic.  Thus there is a
doubly stochastic matrix $B$ such that
\begin{equation}\label{eq:domination}
 \lambda B\le A\quad\hbox{entrywise},
 \qquad P\ge\lambda^n\per B.
\end{equation}
Choose $I,J$ attaining \eqref{eq:lambda}, and set
\[
 u=n-|I|,\qquad v=n-|J|,\qquad k=n-u-v>0.
\]
There are three cases: the whole cut $(u,v)=(0,0)$, a marginal cut where
exactly one of $u,v$ is zero, and a proper cut where $u,v\ge1$.

If $u=v=0$, then $\lambda=1$.  Since $B\le A$ and both matrices have total
sum $n$, $A=B$ is doubly stochastic.  The van der Waerden theorem and
$P\le\gamma_n$ force $A=J_n/n$, contrary to the boundary assumption.

\section{The two-independent-zero permanent gap}

Let $\Omega_n^{(2)}$ be the face of the Birkhoff polytope specified by
$b_{12}=b_{21}=0$, and put $m=n-2$.  Minc's averaging theorem permits a
permanent minimizer on this face to be chosen in the form
\begin{equation}\label{eq:Y}
 Y=
 \begin{pmatrix}
 x_1&0&a_1\one^{\mathsf T}\\
 0&x_2&a_2\one^{\mathsf T}\\
 a_1\one&a_2\one&xJ_m
 \end{pmatrix},
 \qquad x_i=1-ma_i,
 \qquad a_1+a_2+mx=1.
\end{equation}
Here $0\le a_i\le1/m$.

\begin{proposition}[One-variable reduction]\label{prop:face-reduction}
For $m\ge4$, the minimum permanent in \eqref{eq:Y} occurs with $a_1=a_2$.
Consequently the minimum of $\per B$ on $\Omega_n^{(2)}$ is the minimum, on
\[
 \frac{m-2}{m^2}\le x\le\frac1m,
\]
of
\begin{equation}\label{eq:Pn}
 P_n(x)=m!x^{m-2}
 \left(x^2a^2+2mxab^2+m(m-1)b^4\right),
\end{equation}
where
\begin{equation}\label{eq:ab}
 a=\frac{2-m+m^2x}{2},\qquad b=\frac{1-mx}{2}.
\end{equation}
\end{proposition}
\begin{proof}
Put $s=a_1+a_2$, $p=a_1a_2$, and $x=(1-s)/m$.  Counting permanent terms in
\eqref{eq:Y} gives, apart from the positive factor $m!x^{m-2}$,
\[
 E_m(s,p)=x^2(1-ms+m^2p)
 +mx\bigl(s^2-(2+ms)p\bigr)+m(m-1)p^2.
\]
For fixed $s$,
\[
 \frac{\partial E_m}{\partial p}
 =2m(m-1)p+(m+1)s^2-ms-1.
\]
Since $p\le s^2/4$ and $0\le s\le2/m$, this is at most
\[
 D_m(s)=\frac{(m^2+m+2)s^2-2ms-2}{2}.
\]
The quadratic is convex, while
\[
 D_m(0)=-1,
 \qquad D_m(2/m)=-1+\frac2m+\frac4{m^2}<0
 \quad(m\ge4).
\]
Thus the derivative is strictly negative.  For fixed $s$, the permanent is
minimized by maximizing $p$, namely by $a_1=a_2=b$.  Solving the stochastic
constraints yields \eqref{eq:ab} and the stated interval.  A second count
gives \eqref{eq:Pn}.
\end{proof}

For $n=8$, set
\begin{equation}\label{eq:L}
 L_8=\frac{12249}{5000000}.
\end{equation}

\begin{proposition}[Exact two-zero bound]\label{prop:two-zero-bound}
Every $B\in\Omega_8^{(2)}$ satisfies
\[
 \per B>L_8>\gamma_8.
\]
\end{proposition}
\begin{proof}
Here $m=6$ and the complete interval is $[1/9,1/6]$.  The primary verifier
expands $P_8(x)-L_8$ over $\QQ$, verifies positive endpoint values, and uses
an exact rational Sturm sequence to prove that it has no root on the interval.

The independent verifier instead evaluates the matrix \eqref{eq:Y} with
\eqref{eq:ab} by the literal exact Ryser formula at nine rational points,
interpolates the degree-eight polynomial, and proves positivity by exact
Bernstein coefficients on 256 subintervals.  Finally,
\begin{equation}\label{eq:L-gap}
 L_8-\gamma_8=
 \frac{476577}{10240000000}>0.
\end{equation}
\end{proof}

\section{Marginal tight cuts}\label{sec:marginal}

We first record the stationarity equation used twice below.

\begin{lemma}[Support stationarity]\label{lem:stationarity}
At a global maximizer, for every positive entry $a_{ij}$,
\begin{equation}\label{eq:stationarity}
 \frac{R}{r_i}+\frac{C}{c_j}-\per A(i\mid j)=\Phi(A),
\end{equation}
where $A(i\mid j)$ is obtained by deleting row $i$ and column $j$.
\end{lemma}
\begin{proof}
The left side is $\partial\Phi/\partial a_{ij}$.  Moving a small amount of
mass between positive entries shows that all partial derivatives on the
support have a common value $\mu$.  Euler's identity, using homogeneity of
degree $n$, gives
\[
 n\Phi(A)=\sum_{i,j}a_{ij}\frac{\partial\Phi}{\partial a_{ij}}
 =\mu\sum_{i,j}a_{ij}=n\mu.
\]
Thus $\mu=\Phi(A)$.
\end{proof}

Suppose the tight cut is marginal.  Its ratio in \eqref{eq:lambda} is an
average of a nonempty set of row sums or column sums.  Singleton marginal cuts
are also allowed, so a minimizing marginal cut equals a minimum line sum.  By
transposition, choose a minimum row $i$ with
\[
 r_i=\lambda=1-a,\qquad 0\le a<1.
\]
Multiply \eqref{eq:stationarity} by $a_{ij}$ and sum over $j$.  Expansion of
the permanent along row $i$ gives
\[
 R+C\sum_j\frac{a_{ij}}{c_j}-P=r_i\Phi(A).
\]
Therefore
\begin{equation}\label{eq:qi}
 q_i:=\sum_j\frac{a_{ij}}{c_j}
 =1-(1+\theta)a,
 \qquad \theta=\frac{R-P}{C}.
\end{equation}
By \eqref{eq:product-lower} and $P=\gamma_n-\delta$,
\begin{equation}\label{eq:theta}
 \theta\ge R-P\ge1-\gamma_n.
\end{equation}
Let $c_{\max}=\max_jc_j$.  Since $q_i>0$,
\[
 \frac{1-a}{c_{\max}}\le q_i\le1-(2-\gamma_n)a.
\]
Thus the last denominator is positive and
\begin{equation}\label{eq:h}
 c_{\max}\ge1+h(a),
 \qquad h(a)=\frac{(1-\gamma_n)a}{1-(2-\gamma_n)a}.
\end{equation}

AM--GM on the other seven row sums and column sums gives
\begin{align}
 R&\le f(a):=(1-a)\left(1+\frac a7\right)^7,
 \label{eq:f}\\
 C&\le g(h(a)):=(1+h(a))\left(1-\frac{h(a)}7\right)^7.
 \label{eq:g}
\end{align}
Define
\begin{equation}\label{eq:D}
 D(a)=2-f(a)-g(h(a)).
\end{equation}
The shared deficit implies $\delta\ge D(a)$.  The functions $f$ and $g$ are
strictly decreasing on the relevant positive arguments, while $h$ is
increasing, so $D$ is strictly increasing.  The exact check
\begin{equation}\label{eq:D-endpoint}
 D(1/20)>\gamma_8
\end{equation}
confines every feasible value to $0\le a<1/20$.

The dominated doubly stochastic matrix $B$ inherits the two independent
zeroes of $A$, so Proposition~\ref{prop:two-zero-bound} gives
\begin{equation}\label{eq:marginal-permanent}
 P\ge L_8(1-a)^8.
\end{equation}
The elementary scalar contradiction is
\begin{equation}\label{eq:H}
 H(a):=L_8(1-a)^8-\gamma_8+D(a)>0.
\end{equation}
After clearing the positive denominator
$7^7[1-(2-\gamma_8)a]^8$, its numerator has degree 16.  The primary verifier
uses exact Sturm sequences to prove $H>0$ on
\begin{equation}\label{eq:safe-intervals}
 [0,1/400]\ \cup\ [1/75,1/20].
\end{equation}
The independent verifier reconstructs the numerator from exact values and
proves the same statements by Bernstein positivity.  It remains to treat
\begin{equation}\label{eq:risk-interval}
 \frac1{400}\le a\le\frac1{75}.
\end{equation}

\subsection{The large-column support bound}

The weighted-average identity
\[
 \frac{q_i}{r_i}=\sum_{j:a_{ij}>0}\frac{a_{ij}}{r_i}\frac1{c_j}
\]
shows that some $j_*$ with $a_{ij_*}>0$ satisfies
\begin{equation}\label{eq:distinguished-column}
 c_{j_*}\ge\frac{r_i}{q_i}\ge1+h(a).
\end{equation}
Thus the large column may be chosen to contain the minimum row.

Put
\begin{equation}\label{eq:q-support}
 d(a)=\left(1-\frac{h(a)}7\right)^7,
 \qquad q(a)=2-\gamma_8-d(a).
\end{equation}
For every $\ell$ with $a_{\ell j_*}>0$, cofactor nonnegativity and
Lemma~\ref{lem:stationarity} give
\[
 \frac{R}{r_\ell}
 =\Phi(A)-\frac{C}{c_{j_*}}+\per A(\ell\mid j_*)
 \ge2-\gamma_8-d(a)=q(a).
\]
Indeed, $C/c_{j_*}$ is the product of the other seven column sums, whose sum
is at most $7-h(a)$.  On \eqref{eq:risk-interval}, $q$ is increasing and the
verifier checks exactly that $q(1/400)>1$.  Since $R\le1$,
\begin{equation}\label{eq:row-cap}
 r_\ell\le\frac{R}{q(a)}\le\frac1{q(a)}<1
 \qquad(a_{\ell j_*}>0).
\end{equation}

Let $z$ be the number of zero entries in column $j_*$.  Full
indecomposability excludes a column with only one positive entry, while
\eqref{eq:row-cap} excludes a full-support column.  Hence $1\le z\le6$.
The support consists of the distinguished row of sum $1-a$ and $7-z$ other
rows bounded by $1/q(a)$.  Optimizing the remaining $z$ row sums by AM--GM
gives
\begin{equation}\label{eq:Uz}
 R\le U_z(a):=(1-a)q(a)^{-(7-z)}
 \left(\frac{7+a-(7-z)/q(a)}{z}\right)^z.
\end{equation}
The objective is increasing in each capped row sum: after the uncapped rows
are equalized, their mean is greater than $1$, whereas every capped row is
less than $1$.

The matrix $B$ also inherits the $z\times1$ zero rectangle in column $j_*$.
Hwang's staircase theorem therefore gives
\begin{equation}\label{eq:Mz}
 \per B\ge M_z:=
 \max\left\{L_8,\frac{\gamma_{8-z}\gamma_7}{\gamma_{7-z}}\right\}.
\end{equation}
The six exact values used by the verifier are
\[
 \left(
 \frac{12249}{5000000},
 \frac{15625}{6353046},
 \frac{36864}{14706125},
 \frac{1215}{470596},
 \frac{320}{117649},
 \frac{360}{117649}
 \right).
\]
Consequently
\begin{equation}\label{eq:R0}
 R\ge R_{0,z}(a):=2-\gamma_8+M_z(1-a)^8-g(h(a)).
\end{equation}

Here is the exact polynomial certificate for
$R_{0,z}>U_z$.  Put
\begin{align*}
 e&=1-(2-\gamma_8)a,
 &t&=e-\frac{(1-\gamma_8)a}{7},\\
 Q&=(2-\gamma_8)e^7-t^7=e^7q(a),\\
 S_z&=(2-\gamma_8+M_z(1-a)^8)e^8
       -(e+(1-\gamma_8)a)t^7=e^8R_{0,z}(a),\\
 W_z&=(7+a)Q-(7-z)e^7.
\end{align*}
All denominators are positive on \eqref{eq:risk-interval}, and the desired
inequality is equivalent to
\begin{equation}\label{eq:Gz}
 G_z(a):=z^zQ^7S_z
 -(1-a)e^{8+7(7-z)}W_z^z>0.
\end{equation}
Each $G_z$ has degree 65.  The primary verifier expands it directly and
converts it to the degree-65 Bernstein basis on $[1/400,1/75]$; every one of
the $6\cdot66=396$ coefficients is strictly positive.  The smallest is
approximately $7.1630735\cdot10^{-6}$.  The independent audit reconstructs
each $G_z$ from 66 exact values and obtains the same positivity conclusion.
This contradicts \eqref{eq:Uz}--\eqref{eq:R0} for every possible support and
excludes all marginal cuts.

\section{Proper tight cuts}

Suppose now that $u,v\ge1$.  Define
\[
 \varepsilon_r=\sum_{i\in I^c}r_i-u,
 \qquad
 \varepsilon_c=\sum_{j\in J^c}c_j-v.
\]
For positive numbers $x_1,\dots,x_n$ of sum $n$ and product $X$, binary
Pinsker after grouping a subset and its complement gives
\[
 \left(\sum_{i\in S}x_i-|S|\right)^2
 \le-\frac n2\log X.
\]
Applying this to the row and column sums, then using Cauchy--Schwarz and
\eqref{eq:product-lower}, yields
\begin{equation}\label{eq:T}
 |\varepsilon_r|+|\varepsilon_c|
 \le\sqrt{-n\log(RC)}
 \le\sqrt{\frac{n\delta}{1-\delta}}=:T.
\end{equation}
The cut identity
\[
 s(A[I,J])=k-\varepsilon_r-\varepsilon_c+s(A[I^c,J^c])
\]
therefore implies
\begin{equation}\label{eq:lambda-proper}
 \lambda\ge1-\frac Tk.
\end{equation}
For $n=8$, the exact inequality $n\gamma_n<1-\gamma_n$ ensures $T<1\le k$.

Because $B$ is doubly stochastic,
\[
 s(B[I,J])=k+s(B[I^c,J^c])\ge k.
\]
Tightness and \eqref{eq:domination} give the reverse bound after
multiplication by $\lambda$, so equality holds and $B[I^c,J^c]=0$.  Thus $B$
lies on the staircase face with a prescribed $u\times v$ zero block.
Hwang's theorem gives
\begin{equation}\label{eq:nu}
 \per B\ge\nu_{n;u,v}
 :=\frac{\gamma_{n-u}\gamma_{n-v}}{\gamma_k},
 \qquad k=n-u-v,
\end{equation}
where $\gamma_0=1$.  The barycenter has
$(n-v)!(n-u)!/k!$ supported permutations, all of the same weight; both
verifiers reconstruct this count exactly.

Using \eqref{eq:lambda-proper}, Bernoulli's inequality, and
$1-\delta\ge1-\gamma_n$ gives
\begin{align*}
 \nu_{n;u,v}\left(1-\frac Tk\right)^n-(\gamma_n-\delta)
 &\ge \nu_{n;u,v}-\gamma_n+\delta
 -\frac{n\nu_{n;u,v}}k
  \sqrt{\frac{n\delta}{1-\gamma_n}}.
\end{align*}
Completing the square in $\sqrt\delta$ gives the uniform lower bound
\begin{equation}\label{eq:eta}
 \eta_{n;u,v}
 :=\nu_{n;u,v}-\gamma_n
 -\frac{n^3\nu_{n;u,v}^2}{4k^2(1-\gamma_n)}.
\end{equation}
Both verifiers check $\eta_{8;u,v}>0$ for 15 of the 21 proper cuts
\[
 u,v\ge1,\qquad u+v\le7.
\]
The smallest positive value occurs at $(u,v,k)=(1,1,6)$, where
\[
 \nu_{8;1,1}=\frac{33592320}{13841287201},
 \qquad
 \eta_{8;1,1}=
 \frac{8906343086971146850328025}
 {3283430548652519934896577839104}>0.
\]
The only cuts not settled by \eqref{eq:eta} have $k=1$, namely
$(u,v)=(1,6),(2,5),(3,4)$ and their transposes.  We treat them with a
sharper product estimate.

For $1\le d\le7$, define
\begin{equation}\label{eq:Fd}
 F_d(x)=\left(1+\frac{x}{d}\right)^d
        \left(1-\frac{x}{8-d}\right)^{8-d}.
\end{equation}
If a $d$-element subset of the row sums has sum $d+x$, AM--GM separately
within the subset and its complement gives $R\le F_d(x)$; the analogous
statement holds for $C$.  Put
\begin{equation}\label{eq:cd}
 (c_1,c_2,c_3,c_4,c_5,c_6,c_7)
 =\left(\frac12,\frac3{10},\frac14,\frac6{25},
        \frac14,\frac3{10},\frac12\right).
\end{equation}
The exact polynomial inequalities
\begin{equation}\label{eq:subset-deficit}
 1-F_d(x)\ge c_dx^2
 \qquad\left(-\frac1{10}\le x\le\frac1{10},\quad1\le d\le7\right)
\end{equation}
are checked as follows.  After removing the double zero at $x=0$, each
$1-F_d(x)-c_dx^2$ has a degree-six quotient.  Both verifiers convert that
quotient to the Bernstein basis on $[-1/10,1/10]$ and find every coefficient
strictly positive; the smallest coefficient is
\[
 \frac{640006399}{65536000000}>0.
\]
They also check $1-F_d(\pm1/10)>\gamma_8$ for every $d$.  Since $F_d$ is
strictly increasing to $1$ on the negative side of zero and strictly
decreasing from $1$ on the positive side, the lower bounds
$R,C\ge1-\delta\ge1-\gamma_8$ localize every relevant subset deviation to
$(-1/10,1/10)$.

Now take one of the six remaining cuts and write
\[
 x=\varepsilon_r,\qquad y=\varepsilon_c,
 \qquad s=x+y,\qquad u+v=7.
\]
The cut identity gives $\lambda\ge1-s$.  If $s\le0$, then $\lambda\ge1$;
as $\lambda\le1$, the domination argument makes $A$ doubly stochastic and
contradicts the boundary assumption.  Hence $s>0$.  From
\eqref{eq:shared-deficit}, \eqref{eq:Fd}, and
\eqref{eq:subset-deficit},
\begin{equation}\label{eq:s-deficit}
 \delta\ge\rho+\sigma
 \ge c_ux^2+c_vy^2
 \ge \frac{c_uc_v}{c_u+c_v}s^2=:c_{u,v}s^2.
\end{equation}
For the three cases up to transposition,
\begin{equation}\label{eq:ceff}
 c_{1,6}=\frac3{16},\qquad
 c_{2,5}=\frac3{22},\qquad
 c_{3,4}=\frac6{49}.
\end{equation}
The weakest of these constants and the exact inequality
$2401\gamma_8<6$ imply $s<1/7$.

Here the staircase core has size one, so \eqref{eq:nu} becomes
$\nu_{8;u,v}=\gamma_{8-u}\gamma_{8-v}$.  Therefore
\[
 P\ge\nu_{8;u,v}(1-s)^8,
 \qquad
 P=\gamma_8-\delta\le\gamma_8-c_{u,v}s^2.
\]
For $(u,v)=(1,6),(2,5),(3,4)$, both verifiers prove
\begin{equation}\label{eq:Kuv}
 K_{u,v}(s):=\nu_{8;u,v}(1-s)^8-\gamma_8+c_{u,v}s^2>0
 \qquad\left(0\le s\le\frac17\right)
\end{equation}
by exact degree-eight Bernstein coefficients on the two intervals
$[0,1/14]$ and $[1/14,1/7]$.  Transposition handles the other three cuts.
The smallest coefficient among these certificates is
\[
 \frac{13183225395}{259172170858496}>0.
\]
Thus \eqref{eq:Kuv} contradicts the preceding upper bound for $P$, and all
proper cuts are excluded.

The whole, marginal, and proper cases exhaust every tight cut.  No boundary
global maximizer exists.  Hwang's positive-support theorem identifies the
sole remaining maximizer as $J_8/8$, proving Theorem~\ref{thm:main}.

\section{Exact computational audit}

No floating-point value is used in a correctness decision.  The primary
standard-library verifier checks:
\begin{enumerate}
\item the bivariate two-zero equalization identity and its sign;
\item the degree-eight permanent polynomial and the lower bound \eqref{eq:L}
by a rational Sturm sequence;
\item denominator and endpoint signs, \eqref{eq:D-endpoint}, and zero Sturm
roots for the degree-16 numerator on both safe intervals;
\item $q(1/400)>1$, all six support bounds $M_z$, and all 396 exact
Bernstein coefficients in \eqref{eq:Gz};
\item all seven subset-deficit inequalities \eqref{eq:subset-deficit}, their
endpoint localization checks, and the three polynomials \eqref{eq:Kuv};
\item the staircase barycenter count for all 21 proper cuts and all 15
positive values in \eqref{eq:eta}.
\end{enumerate}

The independent verifier imports no code from the primary verifier.  It uses
a literal exact Ryser sum for the two-zero matrices and Bernstein positivity
on 256 subintervals.  It reconstructs the marginal numerator from 17 exact
values and each support polynomial from 66 exact scalar values.  It also
reconstructs each subset quotient from seven values and each exceptional-cut
polynomial from nine values, checks off-grid values, and proves positivity by
an independently implemented Bernstein conversion.  Thus none of the
principal polynomial expansions is shared.

The integration test runs both programs in ordinary and optimized Python mode
and requires identical output.  Three mutations inflate $L_8$, omit the
staircase support improvement in the marginal argument, and replace
$c_4=6/25$ by the false value $1/4$; the appropriate independently mutated
program must reject each change.

\section{Scope and status}

This note establishes dimension 8.  Together with the other exact working
proofs in the accompanying repository, the covered dimensions are
\[
 4,5,8,9,10,11,12,13,14,15,16.
\]
Dimensions $6,7$ are not settled by these materials.  The present result
should be described as an exact computer-assisted working proof until the
mathematical reduction and both programs have undergone independent peer
review.

\section*{Reproducibility}

From the \texttt{n8} directory, run
\begin{verbatim}
python3 -I verify_dittert_n8.py
python3 -I audit_dittert_n8_stdlib.py
python3 -I test_n8_package.py
\end{verbatim}
The verifier final lines must be
\begin{verbatim}
CERTIFIED
INDEPENDENT AUDIT CERTIFIED
\end{verbatim}
and the tests must report \texttt{OK}.

\begin{thebibliography}{99}

\bibitem{CheonWanless2012}
G.-S. Cheon and I. M. Wanless,
\emph{Some results towards the Dittert conjecture on permanents},
Linear Algebra Appl. 436 (2012), 791--801.
\href{https://doi.org/10.1016/j.laa.2010.08.041}{doi:10.1016/j.laa.2010.08.041}.

\bibitem{Egorychev1981}
G. P. Egorychev,
\emph{Solution of the van der Waerden problem for permanents},
Soviet Math. Dokl. 23 (1981), 619--622.

\bibitem{Falikman1981}
D. I. Falikman,
\emph{A proof of van der Waerden's conjecture on the permanent of a stochastic matrix},
Mat. Zametki 29 (1981), 931--938.

\bibitem{Hwang1985}
S.-G. Hwang,
\emph{Minimum permanent on faces of staircase type of the polytope of doubly stochastic matrices},
Linear Multilinear Algebra 18 (1985), 271--306.
\href{https://doi.org/10.1080/03081088508817694}{doi:10.1080/03081088508817694}.

\bibitem{Hwang1986}
S.-G. Hwang,
\emph{A note on a conjecture on permanents},
Linear Algebra Appl. 76 (1986), 31--44.
\href{https://doi.org/10.1016/0024-3795(86)90212-0}{doi:10.1016/0024-3795(86)90212-0}.

\bibitem{Li1990}
C.-K. Li,
\emph{On certain convex matrix sets},
Discrete Math. 79 (1990), 323--326.

\bibitem{Minc1984}
H. Minc,
\emph{Minimum permanents of doubly stochastic matrices with prescribed zero entries},
Linear Multilinear Algebra 15 (1984), 225--243.
\href{https://doi.org/10.1080/03081088408817592}{doi:10.1080/03081088408817592}.

\bibitem{PulaSongWanless2011}
K. Pula, S.-Z. Song, and I. M. Wanless,
\emph{Minimum permanents on two faces of the polytope of doubly stochastic matrices},
Linear Algebra Appl. 434 (2011), 232--238.
\href{https://doi.org/10.1016/j.laa.2010.08.022}{doi:10.1016/j.laa.2010.08.022}.

\end{thebibliography}

\end{document}
